% Options for packages loaded elsewhere
\PassOptionsToPackage{unicode}{hyperref}
\PassOptionsToPackage{hyphens}{url}
\documentclass[
]{article}
\usepackage{xcolor}
\usepackage[margin=1in]{geometry}
\usepackage{amsmath,amssymb}
\setcounter{secnumdepth}{-\maxdimen} % remove section numbering
\usepackage{iftex}
\ifPDFTeX
  \usepackage[T1]{fontenc}
  \usepackage[utf8]{inputenc}
  \usepackage{textcomp} % provide euro and other symbols
\else % if luatex or xetex
  \usepackage{unicode-math} % this also loads fontspec
  \defaultfontfeatures{Scale=MatchLowercase}
  \defaultfontfeatures[\rmfamily]{Ligatures=TeX,Scale=1}
\fi
\usepackage{lmodern}
\ifPDFTeX\else
  % xetex/luatex font selection
\fi
% Use upquote if available, for straight quotes in verbatim environments
\IfFileExists{upquote.sty}{\usepackage{upquote}}{}
\IfFileExists{microtype.sty}{% use microtype if available
  \usepackage[]{microtype}
  \UseMicrotypeSet[protrusion]{basicmath} % disable protrusion for tt fonts
}{}
\makeatletter
\@ifundefined{KOMAClassName}{% if non-KOMA class
  \IfFileExists{parskip.sty}{%
    \usepackage{parskip}
  }{% else
    \setlength{\parindent}{0pt}
    \setlength{\parskip}{6pt plus 2pt minus 1pt}}
}{% if KOMA class
  \KOMAoptions{parskip=half}}
\makeatother
\usepackage{graphicx}
\makeatletter
\newsavebox\pandoc@box
\newcommand*\pandocbounded[1]{% scales image to fit in text height/width
  \sbox\pandoc@box{#1}%
  \Gscale@div\@tempa{\textheight}{\dimexpr\ht\pandoc@box+\dp\pandoc@box\relax}%
  \Gscale@div\@tempb{\linewidth}{\wd\pandoc@box}%
  \ifdim\@tempb\p@<\@tempa\p@\let\@tempa\@tempb\fi% select the smaller of both
  \ifdim\@tempa\p@<\p@\scalebox{\@tempa}{\usebox\pandoc@box}%
  \else\usebox{\pandoc@box}%
  \fi%
}
% Set default figure placement to htbp
\def\fps@figure{htbp}
\makeatother
\setlength{\emergencystretch}{3em} % prevent overfull lines
\providecommand{\tightlist}{%
  \setlength{\itemsep}{0pt}\setlength{\parskip}{0pt}}
\usepackage{bookmark}
\IfFileExists{xurl.sty}{\usepackage{xurl}}{} % add URL line breaks if available
\urlstyle{same}
\hypersetup{
  pdftitle={Experimental Design for Automatic Egg Incubators Using a 2\^{}2 factorial Design and RCBD.},
  pdfauthor={STAT.BONVENTURE.M.MIRUKA},
  hidelinks,
  pdfcreator={LaTeX via pandoc}}

\title{Experimental Design for Automatic Egg Incubators Using a 2\^{}2
factorial Design and RCBD.}
\author{STAT.BONVENTURE.M.MIRUKA}
\date{2026-09-23}

\begin{document}
\maketitle

\begin{quote}
OBJECTIVE: To determine whether Temperature, Humidity or their
interaction is the primary factor driving variations in egg hatchability
rate (Y) and quantify individual-level incubator variability.
\end{quote}

\begin{quote}
Experimental design:
\end{quote}

We utilize a 2\^{}2 full Factorial Design in a Randomized Complete Block
Design.

Factor A : Temperature(T) Low(To): standard low (36.5℃) High(T1):
standard optimal(37.8℃)

Factor B : Relative Humidity(H) Low(Ho): standard low(45\% RH) High(H1):
standard optimal(55\% RH)

Blocks : Incubator models - M1, M2, M3, M4

Experimental units: Fertilized eggs.

Replications : r=3runs per model treatment combination.

\begin{quote}
Statistical analysis framework
\end{quote}

Model : Two-way ANOVA with blocking Model equation : Yijk= μ + αi + βj +
(αβ)ij + γk + εijk, where

\begin{verbatim}
                        μ : overall mean hatchability rate
                        αi: main effect of temp level i
                        βj: main effect of humidity level j
                        (αβ)ij: interaction effect
                        γk: blocking effect of model k
                        εijk: random error ~N(0,σ^2)
                        
\end{verbatim}

\begin{quote}
Dataset generation
\end{quote}

\begin{verbatim}
##   incubator_model temp_level humidity_level hatchability_rate
## 1              M1        low            low              62.4
## 2              M1        low            low              60.8
## 3              M1        low            low              63.1
## 4              M1        low           high              68.5
## 5              M1        low           high              67.2
## 6              M1        low           high              69.0
\end{verbatim}

\#\#ANOVA

\#\#\#libraries

\#\#\#Step 1 :Two\_way anova with blocking (RCBD).

\begin{verbatim}
## === TWO_WAY ANOVA TABLE===
\end{verbatim}

\begin{verbatim}
##                           Df Sum Sq Mean Sq F value   Pr(>F)    
## temp_level                 1 2962.6  2962.6 1605.77  < 2e-16 ***
## humidity_level             1 1037.0  1037.0  562.04  < 2e-16 ***
## incubator_model            3 1517.0   505.7  274.08  < 2e-16 ***
## temp_level:humidity_level  1  120.7   120.7   65.39 5.04e-10 ***
## Residuals                 41   75.6     1.8                     
## ---
## Signif. codes:  0 '***' 0.001 '**' 0.01 '*' 0.05 '.' 0.1 ' ' 1
\end{verbatim}

NOTE: by the view of the Anova table p values of extremely less than
0.05 indicate that there is significant effect of the various factors
towards hatchability rates.

\#\#\#Step 2: Calculating effect size (eta-squared) to identify primary
factor.

\begin{verbatim}
## 
## ===EFFECT SIZE (ETA-SQUARED/ PROPORTION OF VARIANCE)===
\end{verbatim}

\begin{verbatim}
##                           Eta_Squared
## temp_level                     0.5186
## humidity_level                 0.1815
## incubator_model                0.2655
## temp_level:humidity_level      0.0211
## Residuals                      0.0132
\end{verbatim}

NOTE: -Temperature levels account for the vast majority of total
variance confirming it as the primary factor. -Humidity effect is
suppressed by temperature as it account for minimal variance on
hatchability rates.

\#\#\#Step 3: Post-hoc analysis for individual incubator models(Tukey
HSD).

\begin{verbatim}
## 
## === POST-HOC TUKEY HSD (Incubator model differences)===
\end{verbatim}

\begin{verbatim}
##   Tukey multiple comparisons of means
##     95% family-wise confidence level
## 
## Fit: aov(formula = hatchability_rate ~ temp_level * humidity_level + incubator_model, data = incubator_data)
## 
## $incubator_model
##             diff        lwr        upr   p adj
## M2-M1  -4.158333  -5.643131  -2.673536 0.0e+00
## M3-M1 -12.050000 -13.534798 -10.565202 0.0e+00
## M4-M1   2.883333   1.398536   4.368131 3.4e-05
## M3-M2  -7.891667  -9.376464  -6.406869 0.0e+00
## M4-M2   7.041667   5.556869   8.526464 0.0e+00
## M4-M3  14.933333  13.448536  16.418131 0.0e+00
\end{verbatim}

NOTE: Incubator model 1 and 4 show a significant difference in
hatchability rates from the rest of the models(2\&3).

\#\#\#Plot Diagnostics.

\pandocbounded{\includegraphics[keepaspectratio]{Incubator_ANOVA-with-RCBD-design_files/figure-latex/unnamed-chunk-6-1.pdf}}
\pandocbounded{\includegraphics[keepaspectratio]{Incubator_ANOVA-with-RCBD-design_files/figure-latex/unnamed-chunk-7-1.pdf}}

\end{document}
